\section{Discussion and limitations}\label{sec:disc}

\paragraph{What the evidence supports.}
Four claims survive the matched-entry protocol. First, budget
consistency is achievable at negligible cost: the capped water-filling
projection held the contract exactly across every deployed step of every
configuration, while the exponent-style formula it replaces violated the
same contract in $97.4\%$ of simulated steps. Second, allocation
direction matters and static priors do not transfer: the pyramid prior
helped neither of our architectures relative to uniform budgets, and its
mirror image destroyed GPT-2 ($+39\%$ perplexity) while leaving Qwen
almost untouched. Third, structural protection of sinks and recency is
non-negotiable---the $84\times$ collapse without sink protection dwarfs
every allocation or scoring choice. Fourth, the decay rate is a real
dial, not a free lunch: it trades perplexity (best at $\gamma=1$) against
late-context retention and multi-instruction robustness (best at
$\gamma<1$), and Proposition~\ref{prop:staleness} explains the mechanism.

\paragraph{What it does not support.}
We report two negative results with the same prominence as the positive
ones. Dispersion-adaptive allocation did not beat uniform allocation on
the GQA model at this scale (16.40 uniform versus 17.02 adaptive at
$\rho=0.2$, although adaptive still dominates the static pyramid by 1.25
points); and default-rate decay cost 1--2 perplexity points on both
models. Both results bound the strength of our claims deliberately: the
case for \laba{} is that it never loses to the static prior on either
architecture and wins clearly on one, not that it dominates everywhere.
Likewise, the needle results are a measured \emph{failure} of the entire
history-only family at $20\%$ budget, ours included; we consider the
retention diagnostic that localises the failure (cache content, not
generation) more useful than a leaderboard in which every score-based
method scores zero at mid depths.

\paragraph{Limitations.}
The study is CPU-scale by construction: models to $0.5$B parameters,
contexts to $4{,}096$ tokens, single-sequence execution at
$\sim$61 tokens/s prefill on the larger model. We do not measure
throughput under batching, kernel fusion, or paged allocators, and our
overheads are wall-clock on a 2-core CPU rather than profiled GPU
kernels; the systems interactions that dominate at production scale
\cite{kwon2023paged,mao2026infra} are outside what this hardware can
measure. Perplexity uses $10$--$32$ evaluation windows (dictated by
compute), which suffices for the effect sizes reported but is smaller
than the full-test-set convention; the needle grid uses two trials per
depth at 2{,}048 and one at 4{,}096 tokens. Our baseline set covers the
eviction family and its layer-allocation axis; we do not compare against
query-aware methods (SnapKV, Quest, DuoAttention) or quantisation (KIVI)
because they change the deployment model rather than the eviction policy,
and the matched-entry protocol we advocate is exactly the instrument
under which such comparisons should be run once hardware permits. The
earlier manuscript's synthetic records are fully superseded: no number in
this paper descends from them.

\paragraph{Reproducibility.}
The artifact contains the engine, all policies, the projector with its
assertion harness, every result JSON and accounting log behind every
table and figure, the scripts that regenerate them, and a GPU replication
protocol that reuses the identical code paths on larger models and
contexts. A single command re-derives each table from the logs; a single
flag re-runs any configuration with a different seed. We believe this is
the appropriate response to the review context that motivated this work:
claims about memory contracts should be checkable at the level of
individual compression steps, not only through downstream benchmarks.
